\section{The algebra and phases of long words}\label{sec:ring-phases}

The ring used below has two relevant features.  Its positive-degree
matrices are nilpotent, which will give torsion, and its long nonzero
words carry compatible positions modulo arbitrarily large odd integers,
which will give detection by finite quotients.  The first feature is an
input from the companion paper.  We derive the second from the word
hierarchy in that paper.

An ideal $I$ of a unital associative algebra is \emph{matrix-nil} if
every square matrix of every finite size with entries in $I$ is
nilpotent.  All algebras in this paper are over $\F_2$.

\begin{proposition}[The companion algebra]\label{prop:ring-input}
There is an infinite finitely presented unital associative algebra
\[
 R=\bigoplus_{d\geq0}R_d=\F_2\oplus I,
 \qquad R_0=\F_2,
 \qquad I=\bigoplus_{d\geq1}R_d,
\]
such that every $R_d$ is finite-dimensional and $I$ is matrix-nil.
This algebra has the word presentation and hierarchy specified in
Proposition~\ref{prop:hierarchy-input}.
\end{proposition}

This is the algebra constructed in \cite[Theorem~3.7]{Companion}.
We now give the precise part of its structure that we use.  In a
presentation by words, a \emph{replacement} is a relation $v=v'$
between nonempty words of equal length, usable in either direction
inside any word.  A \emph{zero declaration} is a relation $v=0$ for a
nonempty word $v$.  The length of a replacement is the common length
of its two sides.  For a word $w$, write $|w|$ for its length and $[w]$
for its image in the relevant algebra.

\begin{proposition}[The companion hierarchy]\label{prop:hierarchy-input}
The algebra $R$ in Proposition~\ref{prop:ring-input} is the first
algebra $A_0$ of a sequence $(A_j)_{j\geq0}$ with the following
properties.  There are positive integers $c,k,h_0,D$, with
$D>10(c+1)(k+1)$, and
\[
 h_j=h_0D^j,\qquad L_j=2\cdot2^{h_j}+1.
\]
For each $j\geq0$:
\begin{enumerate}[label=\textup{(\roman*)},leftmargin=*]
\item
$A_j$ is presented by a finite alphabet $\Sigma_j$, finitely many
replacements, and finitely many zero declarations, all of length at
most $k$.  It is graded by word length.
\item
There is a substitution $E_j$ sending each letter of $\Sigma_{j+1}$
to a word of length $L_j$ over $\Sigma_j$.  Extended by concatenation,
it induces a unital algebra homomorphism $A_{j+1}\to A_j$.
\item
A specified subset of $\Sigma_j$ gives cuts immediately before its
letters.  There are integers $B_j,T_j\geq0$ and a finite set of words
$\mathcal T_j$ such that every nonzero word $w$ of length at least
$T_j$ has the unique decomposition
\begin{equation}\label{eq:token-decomposition}
 w=p_w\,t_1\cdots t_r\,s_w,
 \qquad |p_w|,|s_w|\leq B_j,\qquad t_i\in\mathcal T_j.
\end{equation}
Here $p_w$ precedes the first cut, $s_w$ begins at the last cut, and
the $t_i$ are exactly the intervals between consecutive cuts.  A cut
is counted only if its following letter occurs in the word.
For each $t\in\mathcal T_j$ there is a nonempty word $u_t$ over
$\Sigma_{j+1}$ such that
\begin{equation}\label{eq:token-upper-word}
 1\leq |u_t|\leq k,\qquad
 |t|=L_j|u_t|,\qquad [t]=[E_j(u_t)]\quad\hbox{in }A_j.
\end{equation}
\end{enumerate}
\end{proposition}

These are the applicable conclusions of
\cite[Proposition~2.2 and the choice of parameters in Section~3]{Companion}.
In particular, the cuts in (iii) are recognized from individual
letters; restricting to a subword cannot create or move a cut.  We
will not need arbitrary concatenations of tokens to be nonzero.

We recall the elementary word-basis argument
\cite[Lemma~2.1]{Companion}, since it will also identify exactly which
equalities of long words must preserve their phases.  Form a graph
whose vertices are all words, including the empty word, and whose
edges are replacements in arbitrary contexts.  Quotienting the free
algebra by the replacement relations leaves one basis vector for each
connected component: the defining ideal is precisely the span of
differences of words in a common component.  Imposing the zero
declarations then kills exactly the components containing a word with
a zero-declared subword.  Thus, in every $A_j$, the nonzero replacement
classes form a basis, and
\begin{equation}\label{eq:word-zero-criterion}
 [w]=0\quad\Longleftrightarrow\quad
 \text{$w$ can be replaced by a word containing a zero-declared subword.}
\end{equation}
In particular, a subword of a nonzero word is nonzero.  The product of
two basis classes is either zero or the class of their concatenation.
We will often use the same symbol for a nonzero word and its basis
class when it occurs as an algebra entry.

The finite alphabet also shows directly that each homogeneous part
$R_d$ is finite-dimensional.  Write
\[
 R_{\geq d}=\bigoplus_{e\geq d}R_e.
\]
For $d\geq1$ this is a two-sided ideal, and $R/R_{\geq d}$ is a
finite ring.  The remaining task in this Section is to extract
compatible modular positions from the hierarchy.

\begin{lemma}[Phases of long words]\label{lem:phases}
For every integer $d\geq0$, put
\[
 L=L_0L_1\cdots L_d.
\]
There is an integer $N\geq1$ such that every nonzero word $w$ in $R$
of length at least $N$ has a start phase
$\sigma(w)\in\Z/L\Z$ and an end phase
\[
 \tau(w)=\sigma(w)+|w|\in\Z/L\Z
\]
with the following properties:
\begin{enumerate}[label=\textup{(\roman*)},leftmargin=*]
\item Both phases depend only on the basis class of $w$.
\item If $v$ is a contiguous subword of $w$, has length at least $N$,
and starts at offset $b$ from the beginning of $w$, then
\begin{equation}\label{eq:phase-subword}
 \sigma(v)=\sigma(w)+b.
\end{equation}
\end{enumerate}
The integers $L$ so obtained are odd and arbitrarily large.
\end{lemma}

\begin{proof}
Fix $d$.  First we construct phases for literal nonzero words, without
claiming invariance under replacements.  We work backward through the
levels $j=d,d-1,\ldots,0$.  At level $d+1$, use the trivial phase
modulo the empty product $1$, with threshold $1$.  Suppose that phases
with the subword property have been constructed at level $j+1$ modulo
\[
 L'=L_{j+1}\cdots L_d,
\]
above a threshold $n_{j+1}\geq1$.

For every token $t\in\mathcal T_j$, choose one word $u_t$ satisfying
\eqref{eq:token-upper-word}, and use this same choice at every
occurrence of $t$.  For the decomposition
\eqref{eq:token-decomposition}, set
\[
 U(w)=u_{t_1}\cdots u_{t_r}.
\]
Choose
\begin{equation}\label{eq:literal-threshold}
 n_j\geq\max\{1,T_j,\,2B_j+L_j n_{j+1}\}.
\end{equation}
Then every nonzero word $w$ of length at least $n_j$ has such a
decomposition and
\[
 |U(w)|=\frac{|w|-|p_w|-|s_w|}{L_j}\geq n_{j+1}.
\]
Moreover, $U(w)$ is nonzero in $A_{j+1}$.  Indeed, the substitution
homomorphism and the token equalities give
\[
 [w]=[p_w]\,[E_j(U(w))]\,[s_w],
\]
so $[U(w)]=0$ would imply $[w]=0$.  Define the literal start phase by
\begin{equation}\label{eq:phase-induction}
 \sigma_j(w)=L_j\sigma_{j+1}(U(w))-|p_w|
 \pmod{L_jL'}.
\end{equation}
Multiplication by $L_j$ makes this independent of the representative
chosen for the residue $\sigma_{j+1}(U(w))$ modulo $L'$.

We verify the subword property at level $j$.  Let $v$ be a subword of
$w$ of length at least $n_j$, starting at offset $b$.  Its cuts are
exactly the cuts of $w$ whose following letters lie in $v$.
Consequently, its complete tokens are a consecutive
subsequence of the complete tokens of $w$.  Our fixed token choices
make $U(v)$ a literal subword of $U(w)$.  If $a$ is its upper-level
offset, the lower-level distance to its first complete token can be
computed in two ways:
\begin{equation}\label{eq:token-offset}
 b+|p_v|=|p_w|+L_j a.
\end{equation}
Both upper words meet the upper-level threshold.  Applying the
inductive subword property and then \eqref{eq:token-offset} yields
\begin{align*}
 \sigma_j(v)
 &=L_j\bigl(\sigma_{j+1}(U(w))+a\bigr)-|p_v|\\
 &=L_j\sigma_{j+1}(U(w))-|p_w|+b
 =\sigma_j(w)+b
 \pmod{L_jL'}.
\end{align*}
This completes the backward induction.

At level zero we now have literal phases modulo $L$, with the subword
property for lengths at least $n_0$.  Only invariance under
replacements remains.  Choose
\[
 N=2n_0+k.
\]
Consider one replacement inside a nonzero word of length at least
$N$.  It affects at most $k$ letters, so an unchanged prefix or an
unchanged suffix has length at least $n_0$.  That common subword has
the same offset before and after the replacement, because the
replacement preserves length.  Applying the literal subword property
to both words shows that their phases agree.  Every intermediate
word in a nonzero replacement class is nonzero.  Chaining these
equalities proves that the start phase is constant on each class.
Word length is constant on each class as well, so the end phase is
also invariant.  The subword property
already proved continues to hold after increasing the threshold from
$n_0$ to $N$.  Finally, every $L_j$ is an odd integer greater than one,
so the products $L_0\cdots L_d$ are odd and unbounded.
\end{proof}

\begin{corollary}\label{cor:phase-products}
Fix $L,N,\sigma,\tau$ from Lemma~\ref{lem:phases}.  If $u,v$ are
nonzero basis words of lengths at least $N$, then
\[
 \tau(u)\ne\sigma(v)\quad\Longrightarrow\quad uv=0.
\]
If $uv\ne0$, its phases satisfy
$\sigma(uv)=\sigma(u)$ and $\tau(uv)=\tau(v)$.
\end{corollary}

\begin{proof}
If the concatenation is nonzero, apply
\eqref{eq:phase-subword} to its prefix $u$ and its suffix $v$.  This
gives $\sigma(uv)=\sigma(u)$ and
$\sigma(v)=\sigma(uv)+|u|=\tau(u)$; adding $|v|$ gives the end-phase
identity.
\end{proof}
